% paper/sections/abstract.tex

\begin{abstract}
Classical \PID{} controllers for autonomous underwater vehicle (\AUV{})
navigation achieve near-perfect success under nominal fluid conditions,
yet collapse catastrophically when those conditions shift by modest,
ocean-realistic amounts---dropping from $100\%$ to $3\%$ success as
drag and current increase to real-ocean ranges, a $97$-percentage-point
collapse that establishes the practical necessity of physics-aware training.

We present a systematic study of domain-randomisation strategies for
\AUV{} sim-to-real transfer, comparing \emph{Curriculum Domain
Randomisation} (\CDR{}), Uniform \DR{}, Naive \SAC{}, and a \PID{}
baseline across nine experiments ($3$ conditions $\times$ $3$ seeds
$\times$ $10^6$ steps).
\CDR{} is a performance-gated strategy that begins training with narrow
physics-parameter ranges and expands them as the policy's rolling success
rate exceeds an adaptive threshold.

Evaluated zero-shot on a held-out test distribution with elevated drag
($0.60$--$1.20$\,kg/m) and strong currents ($0.40$--$0.80$\,m/s)
never seen during training, \CDR{} achieves
$98.67\% \pm 1.89\%$ transfer success.
\CDR{} further consumes $1.13\%$ less thrust energy per step
than Uniform \DR{}
($0.639 \pm 0.062$ vs.\ $0.646 \pm 0.036$,
$p {=} 0.894$, $d {=} 0.14$),
a meaningful saving for battery-constrained real \AUV{} missions.
Domain randomisation eliminates the $\pm 33.83$\,pp
inter-seed deployment variance of policies trained without
randomisation.

We open-source our \MuJoCo{} \AUV{} environment with nine randomisable
parameters to facilitate reproducible underwater robotics research.
\end{abstract}